\PassOptionsToPackage{unicode}{hyperref}
\PassOptionsToPackage{hyphens}{url}
\documentclass[11pt,a4paper]{article}

% ── Codificação e língua ──────────────────────────────────────────────────────
\usepackage[T1]{fontenc}
\usepackage[utf8]{inputenc}
\usepackage[brazil]{babel}

% ── Tipografia ────────────────────────────────────────────────────────────────
\usepackage{lmodern}
\usepackage{microtype}

% ── Cores ─────────────────────────────────────────────────────────────────────
\usepackage{xcolor}
\definecolor{javaOrange}{HTML}{E76F00}
\definecolor{javaDeep}{HTML}{BF5700}
\definecolor{javaBlue}{HTML}{1565C0}
\definecolor{javaTeal}{HTML}{00695C}
\definecolor{javaAmber}{HTML}{B45309}
\definecolor{javaInk}{HTML}{2B2140}
\definecolor{javaCodeBg}{HTML}{FFF8F0}
\definecolor{javaRule}{HTML}{FFD7B5}
\definecolor{javaShade}{HTML}{FFF3E8}

% ── Caixas e código ───────────────────────────────────────────────────────────
\usepackage{tcolorbox}
\tcbuselibrary{skins,breakable,listings}
\usepackage{listings}
\usepackage{fvextra}

\lstdefinestyle{javastyle}{
  language=Java,
  basicstyle=\ttfamily\small,
  keywordstyle=\color{javaBlue}\bfseries,
  stringstyle=\color{javaDeep},
  commentstyle=\color{javaTeal}\itshape,
  numberstyle=\tiny\color{gray},
  breaklines=true,
  breakatwhitespace=false,
  showstringspaces=false,
  tabsize=4,
}

\newtcblisting{javacode}{
  listing only,
  listing style=javastyle,
  breakable,
  enhanced,
  colback=javaCodeBg,
  colframe=javaDeep,
  leftrule=4pt,
  rightrule=0.4pt,
  toprule=0.4pt,
  bottomrule=0.4pt,
  arc=2pt,
  fontupper=\small,
  top=4pt, bottom=4pt, left=6pt, right=6pt,
}

\newtcbox{\inlinecode}{
  on line,
  colback=javaCodeBg,
  colframe=javaRule,
  boxrule=0.5pt,
  arc=2pt,
  fontupper=\ttfamily\small\color{javaDeep},
  boxsep=1pt, left=2pt, right=2pt, top=1pt, bottom=1pt,
}

% ── Layout ────────────────────────────────────────────────────────────────────
\usepackage[margin=2.4cm, top=2.8cm, bottom=2.8cm]{geometry}

% ── Cabeçalho / rodapé ────────────────────────────────────────────────────────
\usepackage{fancyhdr}
\pagestyle{fancy}
\fancyhf{}
\renewcommand{\headrulewidth}{0.4pt}
\renewcommand{\headrule}{\hbox to\headwidth{\color{javaRule}\leaders\hrule height \headrulewidth\hfill}}
\fancyhead[L]{\small\sffamily\color{javaDeep} Java Programming MOOC}
\fancyhead[R]{\small\sffamily\color{javaDeep} Resoluções · Lição 1}
\fancyfoot[C]{\small\sffamily\color{gray}\thepage}

% ── Seções ────────────────────────────────────────────────────────────────────
\usepackage{titlesec}
\titleformat{\section}
  {\sffamily\Large\bfseries\color{javaOrange}}
  {\thesection}{0.7em}{}
  [\vspace{2pt}{\color{javaRule}\titlerule[1.2pt]}]
\titleformat{\subsection}
  {\sffamily\large\bfseries\color{javaDeep}}
  {\thesubsection}{0.6em}{}
\titleformat{\subsubsection}
  {\sffamily\normalsize\bfseries\color{javaTeal}}
  {\thesubsubsection}{0.5em}{}

% ── Listas ────────────────────────────────────────────────────────────────────
\usepackage{enumitem}
\setlist[itemize]{itemsep=2pt, topsep=4pt, parsep=0pt}
\setlist[enumerate]{itemsep=2pt, topsep=4pt, parsep=0pt}

% ── Tabelas ───────────────────────────────────────────────────────────────────
\usepackage{booktabs}
\usepackage{array}
\usepackage{colortbl}
\usepackage{longtable}

% ── Hiperlinks ────────────────────────────────────────────────────────────────
\usepackage{hyperref}
\hypersetup{
  colorlinks=true,
  linkcolor=javaBlue,
  urlcolor=javaBlue,
  citecolor=javaBlue,
  pdftitle={Programação em Java --- Resoluções --- Lição 1: Primeiros Passos},
  pdfauthor={Java Programming MOOC · Universidade de Helsinque},
  pdflang={pt-BR},
}

% ── Caixa de destaque (dica/atenção) ─────────────────────────────────────────
\newtcolorbox{destaque}[1][Observação]{
  enhanced, breakable,
  colback=javaShade,
  colframe=javaOrange,
  leftrule=4pt, rightrule=0.4pt, toprule=0.4pt, bottomrule=0.4pt,
  arc=2pt,
  fonttitle=\sffamily\bfseries\color{javaOrange},
  title=#1,
  top=4pt, bottom=4pt,
}

% ── Título personalizado ──────────────────────────────────────────────────────
\makeatletter
\newcommand{\subtitle}[1]{\gdef\@subtitle{#1}}
\renewcommand{\maketitle}{%
  \begin{center}
    {\color{javaRule}\rule{\linewidth}{1pt}}\\[6pt]
    {\Huge\sffamily\bfseries\color{javaOrange} \@title}\\[6pt]
    \ifx\@subtitle\undefined\else
      {\Large\sffamily\color{javaDeep} \@subtitle}\\[4pt]
    \fi
    {\small\sffamily\color{gray} \@author}\\[2pt]
    {\small\sffamily\color{gray} \@date}\\[6pt]
    {\color{javaRule}\rule{\linewidth}{1pt}}
  \end{center}
  \vspace{1em}
}
\makeatother

% ─────────────────────────────────────────────────────────────────────────────
\title{Programação em Java — Resoluções}
\subtitle{Lição 1: Primeiros Passos}
\author{Java Programming MOOC · Universidade de Helsinque — tradução para o português}
\date{2026}
% ─────────────────────────────────────────────────────────────────────────────

\begin{document}
\maketitle
\tableofcontents
\vspace{1.5em}

% ─────────────────────────────────────────────────────────────────────────────
\section{Exercício 1 --- Olá, usuário!}

\begin{javacode}
import java.util.Scanner;

public class OlaUsuario {
    public static void main(String[] args) {
        Scanner scanner = new Scanner(System.in);

        System.out.print("Qual é o seu nome? ");
        String nome = scanner.nextLine();

        System.out.print("Qual é a sua idade? ");
        int idade = Integer.valueOf(scanner.nextLine());

        System.out.println("Olá, " + nome + "! Você tem " + idade + " anos.");
    }
}
\end{javacode}

\begin{destaque}[Exemplo de execução]
\begin{verbatim}
Qual é o seu nome? Pedro
Qual é a sua idade? 21
Olá, Pedro! Você tem 21 anos.
\end{verbatim}
\end{destaque}

% ─────────────────────────────────────────────────────────────────────────────
\section{Exercício 2 --- Calculadora simples}

\begin{javacode}
import java.util.Scanner;

public class CalculadoraSimples {
    public static void main(String[] args) {
        Scanner scanner = new Scanner(System.in);

        System.out.print("Primeiro número: ");
        int a = Integer.valueOf(scanner.nextLine());

        System.out.print("Segundo número: ");
        int b = Integer.valueOf(scanner.nextLine());

        double quociente = (double) a / b;

        System.out.println("Soma: " + (a + b));
        System.out.println("Diferença: " + (a - b));
        System.out.println("Produto: " + (a * b));
        System.out.println("Quociente: " + quociente);
    }
}
\end{javacode}

O quociente usa \inlinecode{(double) a / b} para não perder a parte decimal do resultado.

% ─────────────────────────────────────────────────────────────────────────────
\section{Exercício 3 --- Par ou ímpar}

\begin{javacode}
import java.util.Scanner;

public class ParOuImpar {
    public static void main(String[] args) {
        Scanner scanner = new Scanner(System.in);

        System.out.print("Digite um número inteiro: ");
        int numero = Integer.valueOf(scanner.nextLine());

        if (numero % 2 == 0) {
            System.out.println(numero + " é par.");
        } else {
            System.out.println(numero + " é ímpar.");
        }
    }
}
\end{javacode}

% ─────────────────────────────────────────────────────────────────────────────
\section{Exercício 4 --- Classificação de temperatura}

\begin{javacode}
import java.util.Scanner;

public class ClassificacaoTemperatura {
    public static void main(String[] args) {
        Scanner scanner = new Scanner(System.in);

        System.out.print("Digite a temperatura em °C: ");
        double temperatura = Double.valueOf(scanner.nextLine());

        if (temperatura < 15) {
            System.out.println("Frio");
        } else if (temperatura <= 25) {
            System.out.println("Agradável");
        } else {
            System.out.println("Quente");
        }
    }
}
\end{javacode}

% ─────────────────────────────────────────────────────────────────────────────
\section{Exercício 5 --- IMC}

\begin{javacode}
import java.util.Scanner;

public class CalculadoraIMC {
    public static void main(String[] args) {
        Scanner scanner = new Scanner(System.in);

        System.out.print("Peso (kg): ");
        double peso = Double.valueOf(scanner.nextLine());

        System.out.print("Altura (m): ");
        double altura = Double.valueOf(scanner.nextLine());

        double imc = peso / (altura * altura);

        System.out.println("Seu IMC é: " + imc);

        if (imc < 18.5) {
            System.out.println("Classificação: Abaixo do peso");
        } else if (imc < 25) {
            System.out.println("Classificação: Peso normal");
        } else if (imc < 30) {
            System.out.println("Classificação: Sobrepeso");
        } else {
            System.out.println("Classificação: Obesidade");
        }
    }
}
\end{javacode}

\begin{destaque}
Os limites usam \inlinecode{<} (estrito) em vez de \inlinecode{<=} para os intervalos, seguindo exatamente as faixas descritas no enunciado (ex.: 24,9 já cai em ``abaixo de 25'').
\end{destaque}

\end{document}
